{\footnotesize
\begin{longtable}{@{}p{0.4cm}>{\raggedright\arraybackslash}p{6.25cm}>{\raggedright\arraybackslash}p{1.8cm}p{0.7cm}r@{}}
\caption{Indicadores del IICA-I (internacional), año utilizado y valor de México en su escala original. ($-$): menor es mejor. Los conteos se expresan por 100 mil habitantes. Fuentes: \parencite{bancomundial_wdi,dsp_v8,girai2026,github_innovationgraph,undesa_egov2024,vdem_v16,wipo_gii2026,wjp_rolindex2026}.}\label{tab:ind-intl}\\
\toprule
D. & Indicador & Fuente & Año & México\\
\midrule
\endfirsthead
\toprule
D. & Indicador & Fuente & Año & México\\
\midrule
\endhead
\bottomrule
\endfoot
1 & Personas que usan internet (\% de la población) & Banco Mundial & 2024 & 83{,}1\\
1 & Matrícula en educación terciaria (\% bruto) & Banco Mundial & 2023 & 47{,}5\\
1 & Índice del componente deliberativo (v2xdl\_delib) & V-Dem v16 & 2025 & 0{,}43\\
2 & Índice de participación de la sociedad civil (v2x\_cspart) & V-Dem v16 & 2025 & 0{,}49\\
2 & Colaboración con otras economías en GitHub (suma anual) [por 100 mil hab.] & GitHub IG & 2025 & 189{,}1\\
3 & Puntaje GIRAI (global) & GIRAI 2026 & 2026 & 27{,}9\\
3 & Contenido de la protección legal de la privacidad (v2smprivcon) & DSP v8 & 2025 & 0{,}81\\
3 & Índice de Servicios en Línea (OSI) & ONU E-Gov & 2024 & 0{,}76\\
4 & Índice del componente participativo (v2x\_partip) & V-Dem v16 & 2025 & 0{,}65\\
4 & Factor 3: Gobierno abierto & WJP & 2026 & 0{,}54\\
4 & Índice de E-Participación (EPI) & ONU E-Gov & 2024 & 0{,}74\\
5 & Voz y rendición de cuentas (WGI) & Banco Mundial & 2025 & -0{,}28\\
5 & Control de la corrupción (WGI) & Banco Mundial & 2025 & -0{,}94\\
5 & Índice del componente igualitario (v2x\_egal) & V-Dem v16 & 2025 & 0{,}47\\
6 & Global Innovation Index (puntaje) & GII (OMPI) & 2026 & 30{,}6\\
6 & Cuentas de desarrolladores en GitHub (al cierre del 4.º trimestre) [por 100 mil hab.] & GitHub IG & 2025 & 1\,999\\
\end{longtable}
}
